\par\Needspace{6\baselineskip}\medskip
\begin{center}\hypertarget{native-183-article137}{}\pdfbookmark[2]{137.}{article-137}137.\end{center}
Die \emph{zweite} Aufgabe besteht darin, die Lage jener drei Großkreise zueinander zu bestimmen.
Ihre Lage zueinander hängt von der Lage ihrer gegenseitigen Schnittpunkte und von ihren
Neigungen ab. Wollen wir diese Verhältnisse eindeutig mit klaren, allgemeinen Begriffen
erfassen, sodass wir für die verschiedenen Fälle keine besonderen Abbildungen benötigen, so
müssen wir einige Erläuterungen vorausschicken. \emph{Erstens} sind auf jedem Großkreis die
beiden entgegengesetzten \emph{Richtungen} voneinander zu unterscheiden. Dazu betrachten wir die
eine als vorwärts gerichtet oder positiv, die andere als rückwärts gerichtet oder negativ. Da
diese Wahl an sich völlig willkürlich ist, folgen wir der festen Regel, stets die Richtungen von
$A,A',A''$ nach $B,B',B''$ als positiv anzusehen. Wird zum Beispiel der Schnittpunkt des ersten Kreises
mit dem zweiten durch einen positiven Abstand vom Punkt $A$ angegeben, so ist dieser von $A$
aus in Richtung $B$ zu nehmen (wie bei $D''$ in unserer Abbildung). Wäre der Abstand negativ,
so müsste man ihn auf der anderen Seite von $A$ abtragen. \emph{Zweitens} sind auch die beiden
Halbkugeln, in die jeder Großkreis die ganze Kugel teilt, durch passende Bezeichnungen zu
unterscheiden. Als \emph{obere} Halbkugel bezeichnen wir diejenige, die rechts liegt, wenn man
den Großkreis auf der Innenseite der Kugel in Vorwärtsrichtung durchläuft, und als \emph{untere}
die andere. Die obere Seite entspricht also der nördlichen Halbkugel in Bezug auf die Ekliptik
oder den Äquator, die untere der südlichen Halbkugel.

Mit diesen Festlegungen lassen sich \emph{beide} Schnittpunkte zweier Großkreise leicht
voneinander unterscheiden. Im einen Schnittpunkt geht der erste Kreis von der unteren Seite des
zweiten auf dessen obere Seite über. Gleichbedeutend damit ist, dass der zweite von der oberen
Seite des ersten auf dessen untere Seite übergeht. Im anderen Schnittpunkt verhält es sich
umgekehrt. Welche Schnittpunkte wir in unserer Aufgabe wählen, ist an sich völlig willkürlich.
Um auch hier nach einer festen Regel zu verfahren, wählen wir stets diejenigen Schnittpunkte
($D,D',D''$ in Abb. 4), an denen der dritte Kreis $A''B''$ auf die obere Seite des zweiten $A'B'$, der
dritte auf die obere Seite des ersten $AB$ und der zweite auf die obere Seite des ersten
übergeht. Die Lage dieser Schnittpunkte wird durch \hypertarget{native-184}{}ihre Abstände von den Punkten $A'$
und $A''$, $A$ und $A''$, $A$ und $A'$ bestimmt, die wir einfach mit $A'D,A''D,AD',A''D',AD'',A'D''$
bezeichnen. Die gegenseitigen Neigungen der Kreise sind dann die Winkel, die in den jeweiligen
Schnittpunkten $D,D',D''$ zwischen den in Vorwärtsrichtung verlaufenden Teilen der sich schneidenden
Kreise liegen. Diese Neigungen, die stets zwischen $0$ und $180^0$ anzusetzen sind, bezeichnen
wir mit $\epsilon,\epsilon',\epsilon''$. Die Bestimmung dieser neun unbekannten Größen aus den bekannten hängt offenbar
von derselben Aufgabe ab, die wir in Artikel 55 behandelt haben. Damit erhalten wir die
folgenden Gleichungen
\[
\begin{aligned}
3)\quad&\sin\tfrac12\epsilon\sin\tfrac12(A'D+A''D)=\sin\tfrac12(l''-l')\sin\tfrac12(\gamma''+\gamma')\\[.5ex]
4)\quad&\sin\tfrac12\epsilon\cos\tfrac12(A'D+A''D)=\cos\tfrac12(l''-l')\sin\tfrac12(\gamma''-\gamma')\\[.5ex]
5)\quad&\cos\tfrac12\epsilon\sin\tfrac12(A'D-A''D)=\sin\tfrac12(l''-l')\cos\tfrac12(\gamma''+\gamma')\\[.5ex]
6)\quad&\cos\tfrac12\epsilon\cos\tfrac12(A'D-A''D)=\cos\tfrac12(l''-l')\cos\tfrac12(\gamma''-\gamma').
\end{aligned}
\]
Aus den Gleichungen 3 und 4 ergeben sich $\tfrac12(A'D+A''D)$ und $\sin\tfrac12\epsilon$, aus den beiden übrigen $\tfrac12(A'D-A''D)$ und
$\cos\tfrac12\epsilon$. Daraus erhält man $A'D,A''D$ und $\epsilon$. Die Mehrdeutigkeit bei der Bestimmung der Bögen
$\tfrac12(A'D+A''D)$, $\tfrac12(A'D-A''D)$ aus ihren Tangenswerten wird durch die Bedingung beseitigt, dass $\sin\tfrac12\epsilon$ und
$\cos\tfrac12\epsilon$ positiv sein müssen. Die Übereinstimmung zwischen $\sin\tfrac12\epsilon$ und $\cos\tfrac12\epsilon$ dient zur Kontrolle
der gesamten Rechnung.

Die Größen $AD',A''D',\epsilon',AD'',A'D'',\epsilon''$ bestimmt man auf ganz dieselbe Weise. Die acht dafür benötigten Gleichungen
brauchen wir hier nicht abzuschreiben, da sie sich unmittelbar aus den Gleichungen 3---6
ergeben, wenn man die Größen jeweils wie folgt ersetzt
\[
\begin{array}{r c|c|c|c|c|c}
 & A'D & A''D & \epsilon & l''-l' & \gamma'' & \gamma'\\
\text{durch} & AD' & A''D' & \epsilon' & l''-l & \gamma'' & \gamma\\
\text{oder durch} & AD'' & A'D'' & \epsilon'' & l'-l & \gamma' & \gamma
\end{array}
\]

Eine weitere Kontrolle der gesamten Rechnung ergibt sich aus der Beziehung zwischen den Seiten
und Winkeln des sphärischen Dreiecks, das von den Punkten $D,D',D''$ gebildet wird. Daraus erhält
man die folgenden Gleichungen, die für jede Lage dieser Punkte gelten
\[
\frac{\sin(AD'-AD'')}{\sin\epsilon}=\frac{\sin(A'D-A'D'')}{\sin\epsilon'}=\frac{\sin(A''D-A''D')}{\sin\epsilon''}\mathpunct{\cdot}
\]

Wird schließlich als Grundebene anstelle der Ekliptik der Äquator \hypertarget{native-185}{}gewählt, so ändert sich
an der Rechnung nur, dass an die Stelle der heliozentrischen Orte der Erde $A,A',A''$ die Punkte des
Äquators treten, in denen ihn die Kreise $AB,A'B',A''B''$ schneiden. Für $l,l',l''$ setzt man also die
Rektaszensionen dieser Schnittpunkte ein und für $A'D$ den Abstand des Punkts $D$ vom
zweiten Schnittpunkt usw.
